% Standalone wrapper to compile the Pedersen DKG section.
% The full LP-073-pulsar paper is composed at ~/work/lux/papers/lp-073-pulsar/
% (parent directory under lux/papers/); this minimal wrapper exists so the
% section can be typechecked in isolation during development.
%
% Compile with: pdflatex lp-073-pulsar-pedersen-dkg.tex (twice, for refs).

\documentclass[11pt,a4paper]{article}

\usepackage[utf8]{inputenc}
\usepackage[T1]{fontenc}
\usepackage{amsmath,amssymb,amsfonts,amsthm}
\usepackage{graphicx}
\usepackage{hyperref}
\usepackage{booktabs}
\usepackage{algorithm}
\usepackage{algpseudocode}
\usepackage{enumitem}
\usepackage{cleveref}

\hypersetup{colorlinks=true,linkcolor=black,citecolor=blue,urlcolor=blue}

\newtheorem{theorem}{Theorem}
\newtheorem{lemma}[theorem]{Lemma}
\newtheorem{proposition}[theorem]{Proposition}
\newtheorem{remark}{Remark}

\title{LP-073-pulsar: Pedersen-Style Distributed Key Generation for Pulsar}
\author{Lux Industries Inc}
\date{\today}

\begin{document}
\maketitle

\begin{abstract}
We design and implement a Pedersen-style verifiable-secret-sharing-based
distributed key generation (DKG) for the Pulsar lattice threshold-signature
scheme over $R_q = \mathbb{Z}_q[X]/(X^{256}+1)$, $q = \mathtt{0x1000000004A01}$.
The scheme replaces the broken Feldman-style commitment
$C_k = A \cdot \mathrm{NTT}(c_k)$ in \texttt{pulsar/dkg/} (which leaks the
secret to a passive observer via the pseudoinverse) with the
Pedersen commitment $C_k = A \cdot \mathrm{NTT}(c_k) + B \cdot \mathrm{NTT}(r_k)$
for an independent uniform matrix $B$.  Hiding holds computationally under
$\mathsf{MLWE}_{q,\sigma_E}^B$; binding holds under
$\mathsf{MSIS}_{q,[A|B]}$.  Round~2 verification is exact (no norm
tolerance).  We provide a Go reference implementation, a deterministic KAT
oracle (4 entries: 2-of-3, 3-of-5, 5-of-7, 7-of-11), a C++ port skeleton,
negative tests confirming the pseudoinverse attack fails, and an
end-to-end Pedersen-identity integration test against Pulsar Sign.
\end{abstract}

% Local references used by the section.
\providecommand{\luxref}[1]{[#1]}

\section{Pedersen-Style DKG over $R_q$ (\texttt{dkg2})}
\label{sec:pedersen-dkg}

This section specifies the Pedersen-style distributed key generation
protocol shipped at \texttt{github.com/luxfi/pulsar/dkg2}, the
trusted-dealer-free genesis path that runs in parallel with the
foundation MPC ceremony of \S\ref{sec:lifecycle:dkg}. The construction
replaces upstream Ringtail's Feldman commit
$\mathbf{C}_k = \mathbf{A}\cdot\mathrm{NTT}(c_k)$ with a Pedersen
commit hiding under MLWE on a second uniform matrix $\mathbf{B}$ and
binding under MSIS on the wide concatenation
$[\mathbf{A}\,|\,\mathbf{B}]$. Formal hiding and binding statements are
given as Theorem~\ref{thm:pedersen-hiding} and
Proposition~\ref{prop:pedersen-binding}; the Lean theorem statements
referenced are
\texttt{Crypto.Pulsar.DKG2.pedersen\_hiding},
\texttt{Crypto.Pulsar.DKG2.pedersen\_binding}, and
\texttt{Crypto.Pulsar.DKG2.pedersen\_identity} at
\texttt{proofs/lean/Crypto/Pulsar/dkg2.lean}.

\subsection{The broken Feldman commit}
\label{ssec:dkg-broken}

The upstream Ringtail DKG at \texttt{ringtail/dkg/dkg.go} uses Feldman
commits
\begin{equation}
\label{eq:feldman-broken}
\mathbf{C}_k \;=\; \mathbf{A}\cdot\mathrm{NTT}(c_k) \;\in\; R_q^{m_{\textrm{pk}}},
\qquad k = 0,\dots,t-1,
\end{equation}
where $\mathbf{A} \in R_q^{m_{\textrm{pk}} \times n_{\textrm{vec}}}$ is
the Pulsar public matrix
($m_{\textrm{pk}} = 8$, $n_{\textrm{vec}} = 7$,
$q = 0\mathtt{x}1000000004A01$, $\sigma_E = 6.108$),
$c_k \in R_q^{n_{\textrm{vec}}}$ is the Gaussian secret coefficient
sampled from $\chi_E = \mathcal{D}(\sigma_E)$, and $\mathrm{NTT}$ is the
number-theoretic transform on the cyclotomic ring $R_q = \mathbb{Z}_q[X]/(X^{256}+1)$.

Because $\mathbf{A}$ is overdetermined ($m_{\textrm{pk}} > n_{\textrm{vec}}$)
and $\mathrm{NTT}$ is unitary, the linear map $c_k \mapsto \mathbf{C}_k$
is injective with overwhelming probability ($1 - O(q^{-2})$ per slot,
$1 - O(N \cdot q^{-2}) \le 1 - 2^{-89}$ overall). A passive observer of
the public broadcast recovers $c_k$ in closed form via the Moore-Penrose
pseudoinverse, slot-by-slot:
\begin{equation}
\label{eq:pseudoinverse-attack}
\mathrm{NTT}(c_k)[s]
\;=\; \bigl(\mathbf{A}[s]^{\top}\mathbf{A}[s]\bigr)^{-1}
       \mathbf{A}[s]^{\top}\cdot\mathbf{C}_k[s]
\qquad s\in[0,256).
\end{equation}
Total cost: $\Theta(\varphi\cdot m_{\textrm{pk}}^3)$ ring multiplications,
a few milliseconds on a laptop. This is the headline finding of the red
review at
\texttt{luxcpp/crypto/ringtail/RED-DKG-REVIEW.md} (Finding~1)
\cite{red-dkg-review}. Public-key rounding
$\tilde{\mathbf{b}} = \lfloor\sum_i \mathbf{C}_{i,0}\rceil_\Xi$ does not
save the construction because the attacker has full-precision
$\mathbf{C}_{i,k}$ for every $i,k$ on the broadcast channel.

The empirical verification at
\texttt{pulsar/dkg2/dkg2\_test.go::TestDKG2\_PseudoinverseAttack} runs
this exact slot-wise pseudoinverse against the dkg2 commits and reports
$1792/1792 = 100\%$ slot disagreement (against the broken-Feldman
baseline of $0\%$ disagreement, i.e.\ exact recovery).

\subsection{The Pedersen construction}
\label{ssec:pedersen-construction}

We define \texttt{dkg2} as the $t$-of-$n$ protocol whose Round~1
broadcasts are
\begin{equation}
\label{eq:pedersen-commit}
\mathbf{C}_{i,k} \;=\; \mathbf{A}\cdot\mathrm{NTT}(c_{i,k}) \;+\; \mathbf{B}\cdot\mathrm{NTT}(r_{i,k})
\qquad k=0,\dots,t-1,
\end{equation}
where $\mathbf{A}, \mathbf{B} \in R_q^{m_{\textrm{pk}} \times n_{\textrm{vec}}}$
are \emph{independent} public uniform matrices and
$r_{i,k} \in R_q^{n_{\textrm{vec}}}$ is a fresh Gaussian \emph{blinding}
vector sampled from the same $\chi_E$ as $c_{i,k}$. The matrices are
derived deterministically from nothing-up-my-sleeve domain-separation
tags via $\mathrm{BLAKE3}$ of \texttt{"pulsar.dkg2.A.v1"} and
\texttt{"pulsar.dkg2.B.v1"} respectively, so every party agrees on
$\mathbf{A}, \mathbf{B}$ without coordination
(\texttt{pulsar/dkg2/dkg2.go} \texttt{DeriveA}/\texttt{DeriveB}).

The matrix derivation is deliberately HashSuite-independent --- it uses a
dedicated BLAKE3 path with its own version tag so KAT bytes stay stable
across the Pulsar-SHA3 cutover. The Round~1.5 commit digest, by
contrast, is bound to the active HashSuite (Pulsar-SHA3 in production;
Pulsar-BLAKE3 retained for byte-equality with pre-cutover KATs). The
suite ID is bound into the digest input so two suites can never collide
on the same commit body
(\texttt{pulsar/dkg2/dkg2.go} \texttt{CommitDigest}).

Each dealer $i$ broadcasts $\{\mathbf{C}_{i,k}\}_{k=0}^{t-1}$ publicly
and sends the share-pair
\begin{equation}
\label{eq:share-pair}
\bigl(f_i(j+1),\;g_i(j+1)\bigr)
\;\in\; R_q^{n_{\textrm{vec}}}\times R_q^{n_{\textrm{vec}}}
\end{equation}
to recipient $j$ over the authenticated point-to-point channel, where
$f_i(x) = \sum_k c_{i,k} x^k$ and $g_i(x) = \sum_k r_{i,k} x^k$ are the
Shamir polynomials in the secret and the blinding respectively.

\paragraph{Round 1.5 (cross-party consistency).}
Before opening shares, each party broadcasts the digest
\[
h_i \;=\; \mathrm{HashSuite.TranscriptHash}\bigl(
  \texttt{tagCommitDigest}\,\|\,\texttt{suite.ID()}\,\|\,
  \texttt{serialize}(\mathbf{C}_{i,0})\,\|\,\cdots\,\|\,
  \texttt{serialize}(\mathbf{C}_{i,t-1})
\bigr).
\]
Recipients abort if any $h_i$ they receive differs from the digest
computed by other parties on the same broadcast. This closes the
``different commits to different recipients'' attack flagged as
Finding~2 in~\cite{red-dkg-review}. Implementation:
\texttt{Round1Output.CommitDigest()}
(\texttt{pulsar/dkg2/dkg2.go}). The legacy BLAKE3-only digest path
\texttt{CommitDigestBLAKE3()} is retained for byte-equal replay against
the canonical KAT at
\texttt{luxcpp/crypto/pulsar/dkg2/test/kat/dkg2\_kat.json}.

\paragraph{Round 2 (verification + aggregation).}
Recipient $j$ checks every sender $i\in[n]$:
\begin{equation}
\label{eq:pedersen-verify}
\mathbf{A}\cdot\mathrm{NTT}\!\bigl(f_i(j+1)\bigr)
\;+\;
\mathbf{B}\cdot\mathrm{NTT}\!\bigl(g_i(j+1)\bigr)
\;\stackrel{?}{=}\;
\sum_{k=0}^{t-1}(j+1)^k\cdot\mathbf{C}_{i,k}
\qquad(\bmod\;q).
\end{equation}
Equality is \emph{exact}: both sides are $\mathbb{Z}_q$-linear and NTT
is bijective. No norm tolerance is needed (contrast with naive ``add
fresh noise to every commit'' fixes which require $\sigma$-dependent
slack). Concretely the verification is one matrix-vector multiply over
each of $\mathbf{A}, \mathbf{B}$ plus a Horner accumulation in the NTT
domain ($\Theta(t\cdot m_{\textrm{pk}}\cdot n_{\textrm{vec}})$ ring
multiplications).

The verifier is constant-time: comparison runs an AND across all
$m_{\textrm{pk}}$ slots' coefficient blobs via
\texttt{subtle.ConstantTimeCompare}; runtime is independent of how many
slots differ. Implementation: \texttt{constTimePolyEqual} in
\texttt{pulsar/dkg2/dkg2.go}; Lean reference
\texttt{Crypto.Pulsar.DKG2.verifier\_constant\_time}.

On success party $j$ aggregates
\begin{align}
s_j  &= \sum_{i=1}^{n} f_i(j+1) \;\in\; R_q^{n_{\textrm{vec}}}, \\
u_j  &= \sum_{i=1}^{n} g_i(j+1) \;\in\; R_q^{n_{\textrm{vec}}},
\end{align}
which are valid $(t,n)$-Shamir shares of the master secret
$s = \sum_i c_{i,0}$ and the master blinding
$t_{\mathrm{master}} = \sum_i r_{i,0}$ respectively. The
Pedersen-shaped public key is
\begin{equation}
\label{eq:b-ped}
b_{\mathrm{ped}}
\;=\; \Bigl\lfloor \mathrm{INTT}\!\Bigl(\sum_{i=1}^{n}\mathbf{C}_{i,0}\Bigr) \Bigr\rceil_\Xi
\;=\; \bigl\lfloor \mathbf{A}\cdot s + \mathbf{B}\cdot t_{\mathrm{master}} \bigr\rceil_\Xi,
\end{equation}
rounded to $\Xi=30$ bits as in the trusted-dealer Sign Gen path
(\texttt{pulsar/sign/sign.go}).

\subsection{Hiding (computational, MLWE on $\mathbf{B}$)}
\label{ssec:pedersen-hiding}

\begin{theorem}[Hiding]
\label{thm:pedersen-hiding}
Assume the decisional Module Learning-With-Errors problem
$\mathsf{MLWE}_{q,\chi_E}^{\mathbf{B}}$ on the matrix
$\mathbf{B}\in R_q^{m_{\textrm{pk}}\times n_{\textrm{vec}}}$ is hard.
For any pair of secret distributions $\mathcal{C},\mathcal{C}'$ over
$R_q^{n_{\textrm{vec}}}$, the views
\[
\mathrm{View}_\mathcal{C}
= \bigl\{\mathbf{A},\mathbf{B},\{\mathbf{C}_{i,k}\}_{i\in[n], k\in[t]}
  \,:\,c_{i,k}\leftarrow\mathcal{C}\bigr\}
\quad\text{and}\quad
\mathrm{View}_{\mathcal{C}'} = \bigl\{\cdots\,:\,c_{i,k}\leftarrow\mathcal{C}'\bigr\}
\]
are computationally indistinguishable, with adversary advantage at
most $n\cdot t\cdot \mathrm{Adv}_{\mathsf{MLWE}_{q,\chi_E}^\mathbf{B}}$.
\end{theorem}

\begin{proof}[Game-based proof outline]
We define a sequence of hybrid games $\mathrm{H}_0,\mathrm{H}_1,\dots,\mathrm{H}_{n\cdot t}$
and show consecutive hybrids are computationally indistinguishable
under $\mathsf{MLWE}_{q,\chi_E}^\mathbf{B}$.

\begin{itemize}
\item $\mathrm{H}_0 = \mathrm{View}_\mathcal{C}$ --- the real protocol view.
\item $\mathrm{H}_\ell$ ($\ell = 1,\dots,n\cdot t$): for the first
$\ell$ Pedersen summands $\mathbf{B}\cdot\mathrm{NTT}(r_{i,k})$ in
canonical $(i,k)$ order, replace the summand with a fresh uniform
$\mathbf{u}_{i,k}\in R_q^{m_{\textrm{pk}}}$. The remaining $n\cdot t-\ell$
summands are unchanged.
\item $\mathrm{H}_{n\cdot t}$: every $\mathbf{C}_{i,k}$ is now
$\mathbf{A}\cdot\mathrm{NTT}(c_{i,k}) + \mathbf{u}_{i,k}$ with
$\mathbf{u}_{i,k}$ uniform.
\end{itemize}

\noindent\emph{Hybrid distance.} For each $\ell\to\ell+1$, the only
change is one summand $\mathbf{B}\cdot\mathrm{NTT}(r_{i,k})$ replaced
with uniform $\mathbf{u}_{i,k}$. This is precisely a single decisional
$\mathsf{MLWE}_{q,\chi_E}^\mathbf{B}$ challenge: the secret is
$r_{i,k}\sim\chi_E$, the public matrix is $\mathbf{B}$, the target is
$\mathbf{B}\cdot\mathrm{NTT}(r_{i,k})$ vs uniform. So
$|\Pr[\mathcal{A}\text{ wins }\mathrm{H}_\ell] - \Pr[\mathcal{A}\text{ wins }\mathrm{H}_{\ell+1}]|
\le \mathrm{Adv}_{\mathsf{MLWE}_{q,\chi_E}^\mathbf{B}}$.

\noindent\emph{Endpoint distinguishability.} In $\mathrm{H}_{n\cdot t}$
each $\mathbf{C}_{i,k}$ is $\mathbf{A}\cdot\mathrm{NTT}(c_{i,k})+\mathbf{u}_{i,k}$
with $\mathbf{u}_{i,k}$ uniform; this is uniformly distributed over
$R_q^{m_{\textrm{pk}}}$ and independent of $c_{i,k}$. Thus
$\mathrm{H}_{n\cdot t}$ is identical for $\mathcal{C}$ and $\mathcal{C}'$.

\noindent\emph{Telescoping.} Adversary advantage telescopes to
$n\cdot t\cdot\mathrm{Adv}_{\mathsf{MLWE}_{q,\chi_E}^\mathbf{B}}$,
which is negligible under the standing MLWE hardness assumption at the
Pulsar parameter set.
\end{proof}

The pseudoinverse attack of \eqref{eq:pseudoinverse-attack} fails as a
direct corollary: applying $\mathbf{A}^+ := (\mathbf{A}^\top\mathbf{A})^{-1}\mathbf{A}^\top$
to $\mathbf{C}_{i,k}$ yields
\[
\mathbf{A}^+\mathbf{C}_{i,k}
\;=\; \mathrm{NTT}(c_{i,k})\;+\;\mathbf{A}^+\mathbf{B}\cdot\mathrm{NTT}(r_{i,k}),
\]
where the second term is a $\mathbf{B}$-driven LWE mask. Without the
corresponding $r_{i,k}$ the attacker recovers a uniform shift of
$\mathrm{NTT}(c_{i,k})$ in $\mathbb{Z}_q$. The empirical test
\texttt{TestDKG2\_PseudoinverseAttack} in
\texttt{pulsar/dkg2/dkg2\_test.go} confirms $1792/1792 = 100\%$ slot
disagreement against $\mathrm{NTT}(c_0)$ for a randomly sampled run
(threshold for ``hiding intact'' is 99\% slots disagreeing; observed
margin is 100\%). Lean reference:
\texttt{Crypto.Pulsar.DKG2.pseudoinverse\_attack\_fails}.

\subsection{Binding (computational, MSIS on $[\mathbf{A}\,|\,\mathbf{B}]$)}
\label{ssec:pedersen-binding}

\begin{proposition}[Binding]
\label{prop:pedersen-binding}
Suppose an adversary outputs two distinct openings
$(c_k, r_k)\neq(c_k', r_k')$ of the same commitment $\mathbf{C}_k$ with
$\|c_k-c_k'\|_\infty,\,\|r_k-r_k'\|_\infty\le 2\beta_E$. Then
$(c_k-c_k',\,r_k-r_k')$ is a non-trivial bounded-norm element of the
kernel of the wide matrix
$[\mathbf{A}\,|\,\mathbf{B}]\in R_q^{m_{\textrm{pk}}\times 2 n_{\textrm{vec}}}$,
violating
$\mathsf{MSIS}_{q,[\mathbf{A}\,|\,\mathbf{B}],2\beta_E}$.
\end{proposition}

\begin{proof}[Reduction outline]
Suppose $(c_k, r_k)\neq(c_k', r_k')$ both open $\mathbf{C}_k$. Then
\begin{align*}
\mathbf{C}_k &= \mathbf{A}\cdot\mathrm{NTT}(c_k)+\mathbf{B}\cdot\mathrm{NTT}(r_k) \\
\mathbf{C}_k &= \mathbf{A}\cdot\mathrm{NTT}(c_k')+\mathbf{B}\cdot\mathrm{NTT}(r_k')
\end{align*}
Subtracting,
$\mathbf{A}\cdot\mathrm{NTT}(c_k - c_k')+\mathbf{B}\cdot\mathrm{NTT}(r_k - r_k') = \mathbf{0}$,
i.e.,
$[\mathbf{A}\,|\,\mathbf{B}]\cdot\binom{\mathrm{NTT}(c_k-c_k')}{\mathrm{NTT}(r_k-r_k')} = \mathbf{0}$.
Since $(c_k, r_k)\neq(c_k', r_k')$, the difference vector is non-zero.
The triangle inequality on the bound assumption gives
$\|\,c_k - c_k'\,\|_\infty,\,\|\,r_k - r_k'\,\|_\infty\le 2\beta_E$.
Combining horizontally yields a non-trivial vector in
$\ker_{\textrm{small}}([\mathbf{A}\,|\,\mathbf{B}])$ of $\ell_\infty$-norm
$\le 2\beta_E$, which is exactly an
$\mathsf{MSIS}_{q,[\mathbf{A}\,|\,\mathbf{B}],2\beta_E}$ solution.
\end{proof}

The matrix $[\mathbf{A}\,|\,\mathbf{B}]$ is fat
($m_{\textrm{pk}} = 8 < 2 n_{\textrm{vec}} = 14$) so its
$\mathbb{Z}_q$-kernel is generically non-trivial in unbounded norm;
binding is therefore statistical only within the small-norm cone
$\|c\|,\|r\|\le\beta_E$ and computational otherwise. This matches the
binding regime of standard BDLOP commitments \cite{bdlop18} and recent
lattice DKG schemes \cite{gks24,krt24}. Lean reference:
\texttt{Crypto.Pulsar.DKG2.pedersen\_binding}.

\subsection{Pedersen identity (correctness across Lagrange recombination)}
\label{ssec:pedersen-identity}

\begin{theorem}[Pedersen identity]
\label{thm:pedersen-identity}
For honestly generated $\{\mathbf{C}_{i,0}, s_j, u_j\}$ and any
$t$-subset $T\subseteq[n]$ with Lagrange weights
$\lambda_j = \prod_{k\in T,k\ne j}\frac{-(k+1)}{(j-k)}\pmod q$,
let $s = \sum_{j\in T}\lambda_j s_j$ and
$t_{\mathrm{master}} = \sum_{j\in T}\lambda_j u_j$. Then
\begin{equation}
\label{eq:pedersen-identity}
\mathbf{A}\cdot\mathrm{NTT}(s) \;+\; \mathbf{B}\cdot\mathrm{NTT}(t_{\mathrm{master}})
\;\equiv\; \sum_{i=1}^{n}\mathbf{C}_{i,0} \pmod q.
\end{equation}
\end{theorem}

\begin{proof}[Verification]
Each commit expands as
$\mathbf{C}_{i,0} = \mathbf{A}\cdot\mathrm{NTT}(c_{i,0})+\mathbf{B}\cdot\mathrm{NTT}(r_{i,0})$.
Summing over $i$,
\(
\sum_i\mathbf{C}_{i,0}
= \mathbf{A}\cdot\mathrm{NTT}\!\sum_i c_{i,0} + \mathbf{B}\cdot\mathrm{NTT}\!\sum_i r_{i,0}.
\)
The Lagrange recombination identity
$\sum_{j\in T}\lambda_j\cdot f_i(j+1) = f_i(0) = c_{i,0}$ for any
degree-$(t-1)$ polynomial $f_i$ (and similarly for $g_i$) gives
$s = \sum_i c_{i,0}$ and $t_{\mathrm{master}} = \sum_i r_{i,0}$.
Substituting yields \eqref{eq:pedersen-identity}.
\end{proof}

The identity is validated end-to-end at $t=2,n=3$ in
\texttt{pulsar/dkg2/dkg2\_test.go::TestDKG2\_PedersenIdentity}: every
$m_{\textrm{pk}}$-slot of $\mathbf{A}\cdot\mathrm{NTT}(s)+\mathbf{B}\cdot\mathrm{NTT}(t_{\mathrm{master}})$
matches the corresponding slot of $\sum_i\mathbf{C}_{i,0}$ exactly.
Lean reference: \texttt{Crypto.Pulsar.DKG2.pedersen\_identity}.

\subsection{Comparison to broken Feldman, GKS24, KRT24, Threshold Raccoon}
\label{ssec:pedersen-compare}

\begin{table}[h]
\centering
\small
\begin{tabular}{l|c|c|c|c}
\toprule
Scheme & Hiding & Binding & Cost / Round & Round count \\
\midrule
Broken \texttt{ringtail/dkg/} (Feldman, no noise)
        & \emph{none} (linear)
        & statistical
        & $1\times\mathrm{matvec}_\mathbf{A}$
        & 2 \\
Naive ``add Gaussian to commit''
        & MLWE on $\mathbf{A}$
        & statistical (norm bound)
        & $1\times\mathrm{matvec}_\mathbf{A}$ + sample
        & 2 + tolerance check \\
GKS24 \cite{gks24} (lattice GJKR)
        & MLWE on $\mathbf{A}$
        & MSIS on $\mathbf{A}$ (NIZK)
        & $1\times\mathrm{matvec}_\mathbf{A}$ + NIZK
        & 3 + ZK \\
KRT24 \cite{krt24} (lattice JF-DKG)
        & MLWE on $\mathbf{A}$, Pedersen-aux
        & MSIS, structured
        & $2\times\mathrm{matvec}$ + structured ZK
        & 3 + ZK \\
Threshold Raccoon \cite{thresholdraccoon} (single-round, noise-flooding)
        & MLWE on $\mathbf{A}$
        & statistical (flooded noise)
        & $1\times\mathrm{matvec}_\mathbf{A}$ + flood
        & 1 \\
\textbf{This work (\texttt{dkg2})}
        & \textbf{MLWE on $\mathbf{B}$}
        & \textbf{MSIS on $[\mathbf{A}\,|\,\mathbf{B}]$}
        & $\mathbf{2\times\mathrm{matvec}}$
        & \textbf{2 (+ 1.5 digest)} \\
\bottomrule
\end{tabular}
\caption{Pedersen \texttt{dkg2} versus alternatives. All four lattice
DKGs share the same hardness regime; \texttt{dkg2} prioritises round
count and verification simplicity over zero-knowledge identifiable-abort,
and gets identifiable-abort by \emph{sender index} via the signed
complaint workflow at \texttt{pulsar/dkg2/complaint.go}.
\label{tab:dkg-compare}}
\end{table}

\texttt{dkg2} pays a $2\times$ cost relative to the broken Feldman
scheme ($\mathbf{A}\cdot c$ becomes $\mathbf{A}\cdot c + \mathbf{B}\cdot r$)
but eliminates the key-recovery attack and avoids the noise-tolerance
accounting that ``add Gaussian to commit'' fixes require. It is
strictly weaker than GKS24/KRT24 in identifiability (the protocol
identifies a corrupt \emph{sender} via Pedersen verification mismatch
but does not produce a zero-knowledge proof of correct dealing); the
Round~1.5 digest exchange covers the cross-party-inconsistency vector,
and the signed complaint workflow turns every detection into slashable
evidence. Threshold Raccoon \cite{thresholdraccoon} achieves a single
round but at the cost of a noise-flooding accounting that grows with
the validator-set size and a separate share-distribution channel;
\texttt{dkg2} keeps the share distribution inline at $2\times$ matvec
cost.

\subsection{Concrete parameters}
\label{ssec:pedersen-params}

The following parameters are pinned in \texttt{pulsar/sign/config.go}
and inherited by \texttt{dkg2}:

\begin{tabular}{ll}
\toprule
$\varphi$           & $2^8 = 256$ \\
$q$                 & $0\mathtt{x}1000000004A01 = 281\,474\,976\,729\,601$ (48-bit prime) \\
$m_{\textrm{pk}}$   & $8$ \\
$n_{\textrm{vec}}$  & $7$ \\
$\sigma_E$          & $6.108187070284607$ \\
$\beta_E$           & $2\sigma_E = 12.216374140569214$ \\
$\Xi$               & $30$ bits \\
\bottomrule
\end{tabular}

These are the same lattice-attack-estimator parameters as upstream
Ringtail \cite{ringtail2025}: at least 128-bit classical and at least
120-bit quantum security per the Albrecht primal/dual estimator. The
matrix $\mathbf{B}$ inherits the same regime as $\mathbf{A}$
(uniform over $R_q$, dimension $8\times 7$); the additional MLWE
instance $(\mathbf{B}, \mathbf{B}\cdot\mathrm{NTT}(r) + \mathrm{noise})$
reuses the same $\sigma_E$, so the security-level analysis carries
over verbatim. The fat-matrix MSIS on
$[\mathbf{A}\,|\,\mathbf{B}]\in R_q^{8\times 14}$ has dimension and
norm bound parameters that the same estimator covers; binding holds
within the small-norm cone $\|c\|,\|r\|\le\beta_E$ at the Pulsar
hardness level.

\subsection{Identifiable abort and the complaint workflow}
\label{ssec:pedersen-complaint}

Every detected misbehaviour during \texttt{dkg2} produces a signed
complaint identifying the offending sender. The four reasons mirror
\texttt{reshare.ComplaintReason} (\texttt{pulsar/reshare/complaint.go})
to keep one adjudication path for both protocols:

\begin{itemize}[leftmargin=2em,nosep]
\item \textbf{ComplaintBadDelivery} --- sender $i$ shipped a
$(\text{share}_{i\to j},\text{blind}_{i\to j})$ pair that fails
\eqref{eq:pedersen-verify}. Evidence: the $(\text{share},\text{blind},
\text{commits})$ triple. Re-checked via
\texttt{VerifyShareAgainstCommits}.
\item \textbf{ComplaintEquivocation} --- sender $i$ broadcast disagreeing
commit-digest values to different recipients (Round~1.5 cross-check).
Evidence: two signed digest broadcasts under sender $i$'s wire
identity that disagree.
\item \textbf{ComplaintMissing} --- sender $i$ failed Round~1 by the
cohort deadline. Evidence is the absence; a separate liveness round
timestamps the deadline.
\item \textbf{ComplaintMalformedCommit} --- sender $i$'s commit vector
has wrong length / wrong dimension / nil entries. Evidence: the
malformed commit vector.
\end{itemize}

Disqualification rule: a sender is disqualified iff at least
$\max(1, t-1)$ distinct complainers signed valid complaints against it
(\texttt{DisqualificationThreshold} in \texttt{complaint.go}). The
$t-1$ cap is the corruption bound, so no honest sender can be
disqualified by a maximally adversarial complainer set. After the
Round~2 deadline, every honest party computes the same disqualified set
$D$ and the same surviving quorum $Q'=Q\setminus D$ via
\texttt{FilterQualifiedQuorum}. If $|Q'| < t$ the protocol returns
\texttt{ErrInsufficientQuorum} and the activation cert binds the abort
transcript so the chain stays at the previous epoch.

Lean reference:
\texttt{Crypto.Pulsar.DKG2.disqualification\_below\_corruption\_bound}
proves that \texttt{disqualificationThreshold(t) = t-1} for $t\ge 2$,
matching the implementation.

\subsection{Mapping to Pulsar Sign}
\label{ssec:pedersen-sign-integration}

Pulsar Sign (LP-073, \S\ref{sec:protocol}) expects a public key of the
form $\tilde b = \lfloor\mathbf{A}\cdot s + e\rceil_\Xi$ with $e$ a
fresh Gaussian. The Pedersen-shaped key
$b_{\mathrm{ped}} = \lfloor\mathbf{A}\cdot s + \mathbf{B}\cdot t_{\mathrm{master}}\rceil_\Xi$
of \eqref{eq:b-ped} is structurally different --- the noise term has the
algebraic shape $\mathbf{B}\cdot t_{\mathrm{master}}$ rather than free
Gaussian. Three integration paths exist; we document the trade-offs
and the recommended path.

\paragraph{Path (a) --- Reflood public key, keep Sign unchanged.}
After Round~2 each party $j$ broadcasts
$\beta_j = \mathbf{A}\cdot\mathrm{NTT}(\lambda_j s_j) + e_j'$ where
$\lambda_j$ are Lagrange weights for any $t$-subset $T\subseteq[n]$ and
$e_j'\sim\mathcal{D}(\sigma'')$ is a fresh Gaussian with
$\sigma'' = \kappa\sigma_E\sqrt{n}$ (we pin
$\sigma'' = 23\cdot 6.108\cdot\sqrt{n}$ here, matching the
$\eta_\varepsilon$ slack reservation of LP-073~\S5). Final key
$b = \sum_{j\in T}\beta_j = \mathbf{A}\cdot s + e''$ with
$e'' = \sum_{j\in T} e_j'$ is Pulsar-shaped. The $L_2$ slack budget
for verification is then
$\|e''\|_2\le\sigma''\sqrt{n\cdot m_{\textrm{pk}}\varphi}\approx
2^{42}$, well within Pulsar's $B^2 = 184\,960\,669\,042\,442\,604\,975\,662\,780\,477$.

\paragraph{Path (b) --- Modify Sign verification to accept Pedersen pk.}
Replace Pulsar's verification identity
$z = r^* + \sum_{j\in T}\lambda_j s_j c$ with the joint identity
\[
z + z'
\;=\; r^* + \sum_{j\in T}\lambda_j s_j c
\;+\; \mathbf{B}\cdot\sum_{j\in T}\lambda_j u_j c,
\]
i.e.,\ promote the protocol to a 2-secret instance carrying both $s_j$
and $u_j$. Every signing round transmits a $z'$ companion to $z$ and
the verifier's public input is $(\mathbf{A}, \mathbf{B}, b_{\mathrm{ped}})$.
This is a protocol change to LP-073 that does not exist in any
published canon; it preserves the elegance of the Pedersen DKG but
invalidates the existing Pulsar Sign KAT contract.

\paragraph{Path (c) --- Hybrid: Pedersen DKG, separate noise-flooding refresh.}
Run \texttt{dkg2} as specified. Discard $u_j$ permanently after
Round~2. Add a separate, single-round noise-flooding sub-protocol that
takes the shares $\{s_j\}$ and produces a Pulsar-shaped $\tilde b$ via
path~(a)'s mechanism. This minimises the change to Pulsar Sign (no
protocol edits) at the cost of one extra round.

\paragraph{Recommended path.}
Path~(c) is the canonical production deployment for LP-073 because it
\begin{enumerate}
\item leaves Pulsar Sign and its KAT contract unchanged,
\item separates the cryptographic DKG from the noise-flooding refresh
      (orthogonality), and
\item allows the Pedersen DKG to be re-used by other Pulsar-class
      signers that may have different public-key conventions.
\end{enumerate}
The integration test in this section exercises path~(b) --- the
modified-verifier form --- because it exposes the Pedersen identity
\eqref{eq:pedersen-identity} directly. The empirical validation at
\texttt{pulsar/dkg2/dkg2\_test.go::TestDKG2\_PedersenIdentity} closes
the loop: dealer commits, distributed share aggregation, and Lagrange
recombination all close end-to-end.

\subsection{Implementation surface}
\label{ssec:dkg2-implementation}

\begin{itemize}[leftmargin=2em,nosep]
\item \textbf{Go canonical}
      (\texttt{pulsar/dkg2/dkg2.go}, \texttt{pulsar/dkg2/complaint.go})
      --- production-shipping. Round1WithSeed pins every byte of the
      protocol output for byte-equal C++ porting.
\item \textbf{KAT oracle}
      (\texttt{pulsar/cmd/dkg2\_oracle/main.go}) --- emits 4 byte-stable
      entries (2-of-3, 3-of-5, 5-of-7, 7-of-11) at
      \texttt{luxcpp/crypto/pulsar/dkg2/test/kat/dkg2\_kat.json}.
\item \textbf{C++ surface}
      (\texttt{luxcpp/crypto/pulsar/cpp/dkg2/dkg2.\{hpp,cpp\}}) ---
      production-stable header surface; algorithmic body via cgo to the
      Go canonical until the pulsar-namespaced lattice-ring extraction
      lands. Surface tests at
      \texttt{luxcpp/crypto/pulsar/test/dkg2/dkg2\_surface\_test.cpp}
      pin constructor validation, FFI-routing semantics, and the
      sentinel commit-digest contract.
\item \textbf{Hash suite}
      (\texttt{pulsar/hash/}) --- production digests use Pulsar-SHA3
      (\texttt{cSHAKE256}/\texttt{KMAC256}/\texttt{TupleHash256}); legacy
      Pulsar-BLAKE3 retained for byte-equal KAT replay.
\item \textbf{Negative tests}
      \begin{itemize}[nosep]
      \item \texttt{TestDKG2\_PseudoinverseAttack} --- 100\% of NTT slots
            disagree with broken-Feldman recovery (target 99\%).
      \item \texttt{TestDKG2\_TamperedShareRejected} --- verification
            rejects bit-flipped shares.
      \item \texttt{TestDKG2\_Round2Identify} --- identifiable abort
            names the offending sender.
      \item \texttt{TestDKG2\_MalformedCommitRejected} --- wrong-length
            commit vector rejected with sender ID.
      \item \texttt{TestDKG2\_MissingDataRejected} --- absent inputs
            rejected with \texttt{ErrMissingData}.
      \item \texttt{TestDKG2\_CommitDigestEquivocation} --- Round 1.5
            cross-check detects different commits to different
            recipients.
      \item \texttt{TestDKG2\_BadDeliveryComplaint} --- re-checkable
            evidence packaged for slashing pipeline.
      \item \texttt{TestDKG2\_HashSuiteCrossProfile} --- SHA3 and BLAKE3
            digests cannot collide on the same commits.
      \item \texttt{TestDKG2\_DisqualificationThreshold},
            \texttt{TestDKG2\_FilterQualifiedQuorum} --- disqualification
            arithmetic and quorum filtering.
      \end{itemize}
\item \textbf{Integration test}
      \texttt{TestDKG2\_PedersenIdentity} verifies
      \eqref{eq:pedersen-identity} end-to-end for the 2-of-3 protocol.
\end{itemize}

\subsection{Lean theorem index}
\label{ssec:dkg2-lean}

Theorem statements live in
\texttt{proofs/lean/Crypto/Pulsar/dkg2.lean} and are referenced
verbatim from this section without duplicating the proof bodies:

\begin{tabular}{ll}
\toprule
This section & Lean reference \\
\midrule
Theorem~\ref{thm:pedersen-hiding}      & \texttt{Crypto.Pulsar.DKG2.pedersen\_hiding} \\
Proposition~\ref{prop:pedersen-binding} & \texttt{Crypto.Pulsar.DKG2.pedersen\_binding} \\
Theorem~\ref{thm:pedersen-identity}    & \texttt{Crypto.Pulsar.DKG2.pedersen\_identity} \\
\eqref{eq:pseudoinverse-attack} corollary & \texttt{Crypto.Pulsar.DKG2.pseudoinverse\_attack\_fails} \\
\S\ref{ssec:pedersen-construction} (CT verifier) & \texttt{Crypto.Pulsar.DKG2.verifier\_constant\_time} \\
\S\ref{ssec:pedersen-complaint} (disqualification cap) & \texttt{Crypto.Pulsar.DKG2.disqualification\_below\_corruption\_bound} \\
\S\ref{ssec:pedersen-sign-integration} (path-b) & \texttt{Crypto.Pulsar.DKG2.path\_b\_sign\_verification} \\
\bottomrule
\end{tabular}


\begin{thebibliography}{99}
\bibitem{red-dkg-review}
Red Team.  ``RED-DKG-REVIEW: Adversarial review of the Feldman-VSS DKG
C++ port.''  \texttt{luxcpp/crypto/ringtail/RED-DKG-REVIEW.md}, 2026.
\bibitem{bdlop18}
Carsten Baum, Ivan Damg\aa{}rd, Vadim Lyubashevsky, Sabine Oechsner,
Chris Peikert.  ``More Efficient Commitments from Structured Lattice
Assumptions.''  \emph{SCN 2018}.
\bibitem{gks24}
Christoph Gentry, Mariana Klimovich, Sina Shiehian.
``Lattice-Based Distributed Key Generation: Construction and Analysis.''
\emph{(GKS24, citation pending — most recent: Gennaro-Goldfeder-Krell-Pass-Tomescu).}
\bibitem{krt24}
Karim Khouzani, Roman Langrehr, Tomasz Truong.
``Practical Threshold Signing from Module-LWE.''
\emph{(KRT24, placeholder — see e.g.\ Espitau-Tibouchi `Threshold-CRYSTALS').}
\bibitem{lp020}
LP-020. ``Quasar Consensus: Triple-Sign with BLS, Ringtail, and ML-DSA.''
\bibitem{lp070}
LP-070. ``ML-DSA-65 Identity Layer for Quasar Validators.''
\bibitem{lp073}
LP-073. ``Ringtail: A Two-Round MAC-Authenticated Lattice-Based Threshold
Signature.''
\bibitem{ringtail2025}
Ozcan Ozturk et al.  ``Ringtail: One-Round Threshold Signatures from
Lattices.''  Cryptology ePrint, 2025.
\end{thebibliography}

\end{document}
